\chapter{A one-sided law of the iterated logarithm}
\label{chap:pre_lil}

The blueprint's almost-sure rates all take the form
$M_n = O(\sqrt{\langle M\rangle_n \log\log\langle M\rangle_n})$. This is the
\emph{upper} half of the law of the iterated logarithm; the sharp constant is
never used. We therefore only formalize the one-sided statement, which is
considerably easier than the full Hartman--Wintner/Stout LIL.

\textbf{Formalization note.} Mathlib has \emph{no} law of the iterated logarithm.
However, \texttt{Probability/Moments/SubGaussian.lean} already provides
conditionally sub-Gaussian moment-generating-function bounds and an
Azuma--Hoeffding tail inequality for martingale differences
(\texttt{measure\_sum\_ge\_le\_of\_HasCondSubgaussianMGF}), together with the
Borel--Cantelli lemmas (\texttt{Probability/BorelCantelli.lean}). These are
exactly the ingredients of the argument below. The bounded-increment case is
within reach of the existing API; the general finite-variance case needs an
additional truncation step.

\section{Exponential tail bound}
\label{sec:pre_lil_tail}

\begin{lemma}[Exponential supermartingale]\label{lem:lil_exp_supermart}
\lean{AlphaRAR.supermartingale_expProcess}\leanok
Let $(M_n)$ be a martingale with $M_0=0$ and increments bounded by
$|\Delta M_i| \le c$. For $|\theta| \le 1/c$, the process
\[
  Z_n(\theta) := \exp\!\Big(\theta M_n - \theta^2 \langle M\rangle_n\Big)
\]
is a supermartingale, where $\langle M\rangle$ is the predictable quadratic
variation (Definition~\ref{def:pred_qv}).
\end{lemma}

\begin{proof}\leanok
\uses{def:pred_qv}
Using $e^x \le 1 + x + x^2$ for $|x| \le 1$ and $\E[\Delta M_i\mid\filt_{i-1}]=0$,
\[
  \E[e^{\theta\Delta M_i}\mid\filt_{i-1}]
  \le 1 + \theta^2\E[(\Delta M_i)^2\mid\filt_{i-1}]
  \le \exp\!\big(\theta^2(\langle M\rangle_i - \langle M\rangle_{i-1})\big).
\]
Hence $\E[Z_i(\theta)\mid\filt_{i-1}] \le Z_{i-1}(\theta)$.
\end{proof}

\begin{lemma}[Freedman-type inequality]\label{lem:lil_freedman}
\lean{AlphaRAR.measure_exists_ge_le_exp}
\uses{lem:lil_exp_supermart}
Under the hypotheses of Lemma~\ref{lem:lil_exp_supermart}, for all $\lambda > 0$
and $v > 0$,
\[
  \Prob\!\Big(\exists n:\ M_n \ge \lambda \ \text{and}\ \langle M\rangle_n \le v\Big)
  \le \exp\!\Big(-\frac{\lambda^2}{2(v + c\lambda)}\Big).
\]
The formalized statement \texttt{AlphaRAR.measure\_exists\_ge\_le\_exp} is the pre-optimization,
finite-horizon form: for each $n$ and each admissible $0<\theta\le 1/c$,
$\Prob(\exists k\le n:\ M_k\ge\lambda,\ \langle M\rangle_k\le v)\le \exp(-\theta\lambda+\theta^2 v)$,
obtained from Ville's maximal inequality
(\texttt{AlphaRAR.smul\_measure\_sup\_le\_integral\_zero}) applied to the exponential
supermartingale. Optimizing over $\theta$ and letting $n\to\infty$ recovers the displayed bound.
\end{lemma}

\begin{proof}
Optional stopping applied to the supermartingale $Z_n(\theta)$ at the first time
$M_n\ge\lambda$ with $\langle M\rangle_n\le v$ gives
$\Prob(\cdots)\le \exp(-\theta\lambda + \theta^2 v)$; optimize over
$0<\theta\le 1/c$. Equivalently, this is the Azuma--Hoeffding/Freedman bound that
follows from the conditionally sub-Gaussian MGF machinery of
\texttt{Probability/Moments/SubGaussian.lean}.
\end{proof}

\begin{corollary}[Optimized Freedman bound]\label{cor:lil_freedman_opt}
\lean{AlphaRAR.measure_exists_ge_le_exp_optimized, AlphaRAR.measure_exists_ge_le_exp_all}\leanok
\uses{lem:lil_freedman}
Under the hypotheses of Lemma~\ref{lem:lil_freedman}, for $\lambda, v > 0$ with the
admissibility condition $\lambda c \le 2v$ (so that the optimizer $\theta=\lambda/(2v)$
satisfies $\theta c\le 1$),
\[
  \Prob\!\Big(\exists n:\ M_n \ge \lambda \ \text{and}\ \langle M\rangle_n \le v\Big)
  \le \exp\!\Big(-\frac{\lambda^2}{4v}\Big).
\]
The variance proxy $2v$ (rather than the sharp $v$) comes from the one-step bound
$e^x\le 1+x+x^2$, costing a factor $\sqrt2$ in the eventual constant.
\end{corollary}

\begin{proof}\leanok
\uses{lem:lil_freedman}
Evaluate Lemma~\ref{lem:lil_freedman} at $\theta=\lambda/(2v)$, giving exponent
$-\theta\lambda+\theta^2 v = -\lambda^2/(4v)$. The finite-horizon events increase to the
event over all $n$, so their measures pass to the limit
(\texttt{measure\_exists\_ge\_le\_exp\_all}).
\end{proof}

\begin{lemma}[Borel--Cantelli over blocks]\label{lem:lil_bc}
\lean{AlphaRAR.ae_eventually_forall_lt_of_summable}\leanok
\uses{cor:lil_freedman_opt}
Fix threshold schedules $(v_k)_{k}$, $(\lambda_k)_k$ with $v_k,\lambda_k>0$ and
$\lambda_k c\le 2v_k$, such that $\sum_k \exp(-\lambda_k^2/(4v_k))<\infty$. Then almost
surely only finitely many blocks are ``bad'': for a.e.\ $\omega$, for all large $k$ and
every $n$, $\langle M\rangle_n\le v_k \Rightarrow M_n < \lambda_k$.
\end{lemma}

\begin{proof}\leanok
\uses{cor:lil_freedman_opt}
By Corollary~\ref{cor:lil_freedman_opt}, $\Prob(s_k)\le\exp(-\lambda_k^2/(4v_k))$ where
$s_k=\{\exists n:\ M_n\ge\lambda_k,\ \langle M\rangle_n\le v_k\}$. The right-hand sides are
summable, so $\sum_k\Prob(s_k)<\infty$, and the first Borel--Cantelli lemma gives that a.e.\
$\omega$ lies in only finitely many $s_k$.
\end{proof}

\section{The upper bound}
\label{sec:pre_lil_upper}

\begin{theorem}[One-sided LIL, bounded increments]\label{thm:lil_bounded}
\uses{lem:lil_bc}
Let $(M_n)$ be a martingale with $M_0=0$ and $|\Delta M_i|\le c$, and suppose
$\langle M\rangle_n \to \infty$ a.s. Then
\[
  \limsup_{n\to\infty}
  \frac{M_n}{\sqrt{2\langle M\rangle_n \log\log\langle M\rangle_n}} \le 1
  \quad\text{a.s.},
\]
and in particular $M_n = O\!\big(\sqrt{\langle M\rangle_n \log\log\langle M\rangle_n}\big)$
a.s.
\end{theorem}

\begin{proof}
\uses{lem:lil_bc}
Fix $\eta>1$ and set $v_k = \eta^k$. Apply Lemma~\ref{lem:lil_bc} with
$\lambda_k=(1+\varepsilon)\sqrt{2v_k\log\log v_k}$: the tail bounds
$\exp(-\lambda_k^2/(4v_k))$ are summable in $k$, so a.s.\ only finitely many blocks are
bad. On the block $\{\langle M\rangle_n\in(v_{k-1},v_k]\}$ this yields
$M_n < \lambda_k \le C\sqrt{\langle M\rangle_n\log\log\langle M\rangle_n}$; letting
$\eta\downarrow1$ and $\varepsilon\downarrow0$ yields the $\limsup$ bound.

\textbf{Formalization status.} The formalized statement
\texttt{AlphaRAR.ae\_eventually\_forall\_lt\_dyadic} is the dyadic-block form: with the schedule
$v_k=2^k$ and $\lambda_k = K\sqrt{v_k(k+1)}$ (where $0<K$ and $Kc\le2$, so admissibility
$\lambda_k c\le 2v_k$ holds for \emph{every} $k$ via $k+1\le 2^k$, and the tail bounds
$\exp(-\lambda_k^2/(4v_k))=\exp(-K^2(k+1)/4)$ form a geometric series), almost surely for all
large $k$ and every $n$, $\langle M\rangle_n\le 2^k \Rightarrow M_n < K\sqrt{2^k(k+1)}$. This is
the complete probabilistic content, built from Ville's inequality, the optimized Freedman bound
(Corollary~\ref{cor:lil_freedman_opt}), and Borel--Cantelli (Lemma~\ref{lem:lil_bc}). Since on the
block $\langle M\rangle_n\in(2^{k-1},2^k]$ one has $\sqrt{2^k(k+1)}\asymp
\sqrt{\langle M\rangle_n\log\langle M\rangle_n}$, it gives the displayed
$O(\sqrt{\langle M\rangle_n\log\langle M\rangle_n})$ upper bound, formalized in
Corollary~\ref{cor:lil_upper} below. One deliberate deviation from the displayed statement remains
unformalized: the $\log\log$ (rather than $\log$) rate with sharp constant $1$, which would require
refining the one-step MGF bound $e^x\le1+x+x^2$ to variance proxy $1$ (valid only as the increment
$\to0$). The $\log$ rate already gives $M_n=o(\langle M\rangle_n)$, which suffices for all
downstream uses.
\end{proof}

\begin{corollary}[Normalized one-sided upper bound]\label{cor:lil_upper}
\uses{thm:lil_bounded}
Under the hypotheses of Theorem~\ref{thm:lil_bounded} (in particular $\langle M\rangle_n\to\infty$
a.s.), there is a deterministic constant $C$ (depending only on $c$) such that almost surely
\[
  M_n \le C\,\sqrt{\langle M\rangle_n\,\log\langle M\rangle_n}\qquad\text{for all large }n.
\]
This is the consumable $O(\sqrt{\langle M\rangle_n\log\langle M\rangle_n})$ a.s.\ upper bound used
downstream (Lemma~\ref{lem:U_increment_bound}, via the bounded assignment martingale).
\end{corollary}

\begin{proof}
\uses{thm:lil_bounded}
Repackage Theorem~\ref{thm:lil_bounded}: for large $n$ let $k$ be the least index with
$\langle M\rangle_n\le 2^k$. Minimality gives $2^k\le 2\langle M\rangle_n$ and
$k\le\log_2\langle M\rangle_n+1$, so
$K\sqrt{2^k(k+1)}\le C\sqrt{\langle M\rangle_n\log\langle M\rangle_n}$ with
$C=K\sqrt{2(1/\log 2+2)}$, using $\log\langle M\rangle_n\ge1$ eventually.
\end{proof}

\begin{corollary}[Upper bound at scale $\sqrt{n\log n}$]\label{cor:lil_upper_nat}
\uses{cor:lil_upper, lem:qv_linear_bound}
Under the hypotheses of Corollary~\ref{cor:lil_upper}, almost surely
$M_n \le C\sqrt{n\log n}$ for all large $n$, for a deterministic constant $C$. This is the form
$M_n = O(\sqrt{n\log n})$ a.s.\ used for the bounded assignment martingale
$M_{\cdot,k}$ in Lemma~\ref{lem:U_increment_bound}.
\end{corollary}

\begin{proof}
\uses{cor:lil_upper, lem:qv_linear_bound}
By Lemma~\ref{lem:qv_linear_bound}, $\langle M\rangle_n\le c^2 n$, hence
$\log\langle M\rangle_n\le\log(c^2 n)\le 2\log n$ once $n\ge c^2$. Substituting into
Corollary~\ref{cor:lil_upper} bounds $\sqrt{\langle M\rangle_n\log\langle M\rangle_n}
\le\sqrt{2c^2}\,\sqrt{n\log n}$.
\end{proof}

\section{The finite-variance case via truncation}
\label{sec:pre_lil_trunc}

The remaining case is a martingale $M_n = \sum_{i<n}\xi_i D_i$ whose increments are
\emph{not} bounded: $\xi_i$ are i.i.d.\ centered with only $\E[\xi_0^2]=\sigma^2<\infty$,
and $D_i\in\{0,1\}$ is predictable (each $\xi_i$ is independent of the past $\filt_{i-1}$,
while $D_i$ is $\filt_{i-1}$-measurable). Then $\langle M\rangle_n=\sigma^2 N_n$ with
$N_n:=\sum_{i<n}D_i$. Theorem~\ref{thm:lil_bounded} does not apply directly, so we
truncate at level $b_i:=\sqrt{i}$ --- the largest level for which the tail is still
summable under only a second moment --- and control three pieces separately. This section
records the decomposition into small steps; each is a candidate Lean lemma to be formalized
in turn. As throughout, we prove the $\log$ (not $\log\log$) rate.

\begin{definition}[Truncated decomposition]\label{def:lil_trunc}
With $b_i:=\sqrt i$, set
$\xi_i^{\mathrm t}:=\xi_i\indic_{\{|\xi_i|\le b_i\}}$,
$\xi_i^{\mathrm r}:=\xi_i\indic_{\{|\xi_i|> b_i\}}$,
$m_i:=\E[\xi_i^{\mathrm t}]$, $\tilde\xi_i:=\xi_i^{\mathrm t}-m_i$, and
\[
  \tilde M_n:=\sum_{i<n}\tilde\xi_i D_i,\quad
  R_n:=\sum_{i<n}\xi_i^{\mathrm r} D_i,\quad
  \mathrm{Dr}_n:=\sum_{i<n} m_i D_i,\quad
  N_n:=\sum_{i<n} D_i,
\]
so that $M_n=\tilde M_n+R_n+\mathrm{Dr}_n$.
\end{definition}

\subsection{The tail term $R_n$}

\begin{lemma}[Summable tail probabilities]\label{lem:trunc_tail_summable}
\lean{AlphaRAR.tsum_measure_abs_sub_gt_sqrt_ne_top}\leanok
\uses{def:lil_trunc}
$\sum_{i}\Prob(|\xi_i|>\sqrt i)<\infty$. (The formalized statement
\texttt{tsum\_measure\_abs\_sub\_gt\_sqrt\_ne\_top} is the law-level version: for any finite measure
$\rho$ on $\R$ with $\int(x-\theta)^2\,d\rho<\infty$, $\sum_i\rho(|x-\theta|>\sqrt i)<\infty$; apply
it to the reward law $\rho=\nu_k$, $\theta=\theta_k$. This rests on a general-measure layer-cake
bound \texttt{tsum\_measure\_Ioi\_ne\_top}, the $\mu$-general form of Mathlib's
\texttt{tsum\_prob\_mem\_Ioi\_lt\_top}.)
\end{lemma}
\begin{proof}\leanok
Since $|\xi_i|>\sqrt i \iff (x-\theta_k)^2>i$, the sum equals $\sum_i\nu_k((\cdot-\theta_k)^2>i)$,
which is finite by the layer-cake bound $\sum_i\mu\{X>i\}\le\int X\,d\mu+\mu(\text{univ})$ applied to
$X=(\cdot-\theta_k)^2$ (integrable, being the arm-$k$ second central moment $\sigma^2$).
\end{proof}

\begin{lemma}[Tail is eventually zero]\label{lem:trunc_tail_const}
\lean{AlphaRAR.ae_eventually_abs_le_of_tsum_ne_top,
AlphaRAR.measure_action_and_tail_le, AlphaRAR.tsum_measure_action_and_tail_ne_top,
AlphaRAR.ae_eventually_tailRespPart_const}\leanok
\uses{lem:trunc_tail_summable}
Almost surely $|\xi_i|\le b_i$ for all large $i$; consequently $\xi_i^{\mathrm r}=0$ eventually and
$R_n$ is a.s.\ eventually constant, so $R_n=O(1)$ a.s.

The formalization has two layers. The abstract eventually-bounded clause
\texttt{ae\_eventually\_abs\_le\_of\_tsum\_ne\_top} is stated for a general level sequence $b$ and
takes $\sum_i\Prob(|\xi_i|>b_i)<\infty$ as hypothesis. The sampled version
\texttt{ae\_eventually\_tailRespPart\_const} concludes that the tail remainder
$R_n=Q_n-\tilde M_n-\mathrm{Dr}_n$ (\texttt{tailRespPart}) is a.s.\ eventually constant, via a
\emph{conditional bookkeeping} step: the sampled tail event has
$\Prob(A_i=k,\ \sqrt i\le|Y_i-\theta_k|)\le\nu_k(\sqrt i\le|\cdot-\theta_k|)$
(\texttt{measure\_action\_and\_tail\_le}, by conditioning on $\filtAction_i$ and
\texttt{condExp\_feedback\_comp}), so these probabilities are summable
(\texttt{tsum\_measure\_action\_and\_tail\_ne\_top}, using the closed-window law-level bound
\texttt{tsum\_measure\_abs\_sub\_ge\_sqrt\_ne\_top}).
\end{lemma}
\begin{proof}\leanok
\uses{lem:trunc_tail_summable}
By Lemma~\ref{lem:trunc_tail_summable} (in the sampled form above) the summability hypothesis holds,
and the first Borel--Cantelli lemma (\texttt{ae\_eventually\_notMem}) gives a.s.\ that for all large
$i$ arm $k$ is not sampled with an off-window response, whence $\xi_i^{\mathrm r}=0$ and the
increments of $R_n$ (\texttt{tailRespPart\_succ\_sub}, where the truncated mean cancels) vanish for
large $i$; a sequence with eventually-vanishing increments is eventually constant.
\end{proof}

\subsection{The drift term $\mathrm{Dr}_n$}

\begin{lemma}[Truncated-mean bound]\label{lem:trunc_mean_bound}
\lean{AlphaRAR.abs_integral_truncation_le}\leanok
\uses{def:lil_trunc}
$|m_i|\le \sigma^2/\sqrt i$ for $i\ge1$. (The formalized statement
\texttt{abs\_integral\_truncation\_le} is the abstract single-variable bound: for a centred $X$
with $X^2$ integrable, $|\E[\operatorname{trunc}(X,A)]|\le \E[X^2]/A$; apply it to $X=\xi_0$ and
$A=\sqrt i$.)
\end{lemma}
\begin{proof}\leanok
Since $\E[\xi_0]=0$, $m_i=\E[\operatorname{trunc}(\xi_0,\sqrt i)]
=\E[\operatorname{trunc}(\xi_0,\sqrt i)-\xi_0]$, so by the pointwise bound
$|\operatorname{trunc}(\xi_0,A)-\xi_0|\le\xi_0^2/A$ (the truncation is either $\xi_0$ or, off the
window where $|\xi_0|\ge A$, is $0$ with $|\xi_0|\le\xi_0^2/A$),
$|m_i|\le\E[\xi_0^2]/\sqrt i=\sigma^2/\sqrt i$.
\end{proof}

\begin{lemma}[Drift is $O(\sqrt{n})$]\label{lem:trunc_drift}
\lean{AlphaRAR.abs_truncDrift_le}\leanok
\uses{lem:trunc_mean_bound}
$|\mathrm{Dr}_n|\le 2\sigma^2\sqrt{n}$ (everywhere, with $\sigma^2=V_k$). In particular
$\mathrm{Dr}_n=O(\sqrt n)=o(\sqrt{N_n\log N_n})$ a.s.\ whenever $N_n$ grows linearly (the case of
interest). The sharper bound $|\mathrm{Dr}_n|\le 2\sigma^2\sqrt{N_n}$ holds by reindexing to the
sampled subsequence $i_j\ge j$, but is not needed and not formalized.
\end{lemma}
\begin{proof}\leanok
\uses{lem:trunc_mean_bound}
By Lemma~\ref{lem:trunc_mean_bound}, $|m_i|\le\sigma^2/\sqrt i$ (and $m_0=0$), so
$|\mathrm{Dr}_n|\le\sum_{i<n}|m_i|D_i\le\sum_{i<n}|m_i|\le\sigma^2\sum_{i<n}\tfrac1{\sqrt i}
\le 2\sigma^2\sqrt n$, using $\sum_{i<n}1/\sqrt i\le 2\sqrt n$
(\texttt{sum\_one\_div\_sqrt\_le}).
\end{proof}

\subsection{The centred truncated martingale $\tilde M_n$}

\begin{lemma}[$\tilde M$ is a martingale]\label{lem:trunc_mart}
\lean{AlphaRAR.martingale_truncRespMart}\leanok
\uses{def:lil_trunc, lem:Q_martingale}
$\tilde M$ is a martingale with $\tilde M_0=0$: $\E[\tilde\xi_i D_i\mid\filt_{i-1}]
=D_i\,\E[\tilde\xi_i\mid\filt_{i-1}]=D_i(\E[\xi_i^{\mathrm t}]-m_i)=0$, using that $\xi_i$ is
independent of $\filt_{i-1}$ and $D_i$ is $\filt_{i-1}$-measurable.
(Formalized in the \texttt{IsAlgEnvSeq}/\texttt{stationaryEnv} framework: $\tilde M$ is the
\emph{truncated response martingale} $\tilde Q_{n,k}=\sum_{i<n}\indic_{\{A_i=k\}}
(\operatorname{trunc}(Y_i-\theta_k,\sqrt i)-m_i)$, a martingale for the action-augmented
filtration $\filtAction$. It is an instance of a general functional response martingale
\texttt{genRespMart} --- of which the response martingale $Q$ (Lemma~\ref{lem:Q_martingale}) is the
$g=\mathrm{id}$ case --- whose martingale property follows from \texttt{condExp\_feedback\_comp}.)
\end{lemma}
\begin{proof}\leanok
\uses{def:lil_trunc}
Pull out the $\filtAction_i$-measurable indicator $\indic_{\{A_i=k\}}$; the conditional
expectation of $g(Y_i)$ (with $g=\operatorname{trunc}(\cdot-\theta_k,\sqrt i)$) given
$\filtAction_i$ is $(\nu(A_i))[g]=m_i$ on $\{A_i=k\}$ (\texttt{condExp\_feedback\_comp}), so the
centred increment has conditional mean $m_i-m_i=0$.
\end{proof}

\begin{lemma}[Quadratic variation of $\tilde M$]\label{lem:trunc_qv}
\lean{AlphaRAR.predQuadVar_truncRespMart_eq}\leanok
\uses{lem:trunc_mart, def:pred_qv, lem:qv_incr}
$\langle\tilde M\rangle_n=\sum_{i<n}v_i D_i$, where $D_i=\indic_{\{A_i=k\}}$ and
$v_i:=\int(\operatorname{trunc}(x-\theta_k,\sqrt i)-m_i)^2\,d\nu_k$ is the truncated-response
variance of arm $k$.
\end{lemma}
\begin{proof}\leanok
\uses{lem:trunc_mart, def:pred_qv, lem:qv_incr}
Instance of the general functional quadratic variation
\texttt{predQuadVar\_genRespMart\_eq}: the compensator increment is the $\filtAction_i$-conditional
second moment $\E[(\tilde\xi_i D_i)^2\mid\filtAction_i]=v_i D_i$
(\texttt{condExp\_genRespMart\_increment\_sq}), and these sum to $\sum_{i<n}v_i D_i$. (Unlike the
response martingale $Q$, whose constant variance gives $V_k N$, the truncated variances $v_i$ vary
with $i$, so the weighted sum is retained.)
\end{proof}

\begin{lemma}[Quadratic variation upper bound]\label{lem:trunc_qv_le}
\lean{AlphaRAR.predQuadVar_truncRespMart_le}\leanok
\uses{lem:trunc_qv}
$\langle\tilde M\rangle_n\le V_k N_{n,k}$ a.s., since each $v_i\le V_k$ (truncation reduces the
absolute value, and variance is at most the second moment). This is the only bound on
$\langle\tilde M\rangle$ used by the block LIL argument (the matching lower bound and the
convergence $v_i\to\sigma^2$ of Lemma~\ref{lem:trunc_var_conv} are not needed for the $O$-rate).
\end{lemma}
\begin{proof}\leanok
\uses{lem:trunc_qv}
By Lemma~\ref{lem:trunc_qv}, $\langle\tilde M\rangle_n=\sum_{i<n}v_i D_i\le V_k\sum_{i<n}D_i
=V_k N_{n,k}$, using $v_i=\operatorname{Var}_{\nu_k}(\operatorname{trunc})\le
\E_{\nu_k}[\operatorname{trunc}^2]\le\E_{\nu_k}[(\cdot-\theta_k)^2]=V_k$.
\end{proof}

\begin{lemma}[Truncated variances converge]\label{lem:trunc_var_conv}
\uses{lem:trunc_qv}
(Not needed for the $O$-rate.) $v_i\to\sigma^2$; in particular $\tfrac12\sigma^2\le v_i\le\sigma^2$
for all large $i$, so
$\tfrac12\sigma^2 N_n\le\langle\tilde M\rangle_n\le\sigma^2 N_n$ eventually and
$\langle\tilde M\rangle_n\asymp\sigma^2 N_n$.
\end{lemma}
\begin{proof}
$v_i=\E[(\xi_i^{\mathrm t})^2]-m_i^2$. As $\sqrt i\to\infty$, $\E[(\xi_i^{\mathrm t})^2]
=\int\operatorname{trunc}(x-\theta_k,\sqrt i)^2\,d\nu_k\to\int(x-\theta_k)^2\,d\nu_k=\sigma^2$
by monotone/dominated convergence, and $m_i\to0$ (Lemma~\ref{lem:trunc_mean_bound}), so
$v_i\to\sigma^2$. The upper bound $v_i\le\sigma^2$ is immediate from
$|\operatorname{trunc}(\xi_i,\sqrt i)|\le|\xi_i|$.
\end{proof}

\subsection{Core: the growing-increment LIL for $\tilde M$}
\label{sec:pre_lil_grow}

The increments $|\tilde\xi_i D_i|\le 2b_i=2\sqrt i$ \emph{grow}, so Theorem~\ref{thm:lil_bounded}
(fixed $c$) is unavailable. We rerun the dyadic argument with a block-dependent increment bound.
On block $k$ we use horizon $2^k$, increment bound $c_k:=2\sqrt{2^k}$, and the \emph{constrained}
exponential parameter $\theta_k:=1/c_k$ (the optimizer is inadmissible). The point is that
$\theta_k^2\langle\tilde M\rangle_{2^k}\le\theta_k^2\sigma^2 2^k=\sigma^2/4$ stays bounded while
$\theta_k\lambda_k$ grows like $\sqrt k$.

\begin{lemma}[Per-block Freedman with constrained $\theta$]\label{lem:trunc_block}
\uses{lem:trunc_mart, lem:lil_freedman}
Fix $C>0$ and set $\lambda_k:=C\sqrt{2^k\,k}$. For each $k$,
\[
  \Prob\!\Big(\exists\,m\le 2^k:\ \tilde M_m\ge\lambda_k,\ \langle\tilde M\rangle_m\le\sigma^2 2^k\Big)
  \le \exp\!\big(-\tfrac{C}{2}\sqrt k+\tfrac{\sigma^2}{4}\big).
\]
Here $\sigma^2=V_k$ is the arm variance and $\langle\tilde M\rangle_m\le V_k 2^k$ uses only the
upper bound Lemma~\ref{lem:trunc_qv_le}.
\end{lemma}
\begin{proof}
\uses{lem:lil_freedman, lem:trunc_qv_le}
Apply the Freedman inequality (Lemma~\ref{lem:lil_freedman}) to the martingale stopped at time
$2^k$ --- whose increments are bounded by $c_k=2\sqrt{2^k}$, so $\theta_k c_k=1$ is admissible ---
with $\theta=\theta_k$, $\lambda=\lambda_k$, $v=\sigma^2 2^k$. The exponent is
$-\theta_k\lambda_k+\theta_k^2\sigma^2 2^k=-\tfrac C2\sqrt k+\tfrac{\sigma^2}4$, since
$\theta_k\lambda_k=\tfrac{C\sqrt{2^k k}}{2\sqrt{2^k}}=\tfrac C2\sqrt k$ and
$\theta_k^2\sigma^2 2^k=\tfrac{\sigma^2 2^k}{4\cdot2^k}=\tfrac{\sigma^2}4$. Stopping at $2^k$ is
needed so the increment bound holds up to the horizon; on the event only times $m\le 2^k$ occur,
where the stopped process agrees with $\tilde M$.

The formalization packages the stopping into a reusable device
\texttt{AlphaRAR.measure\_exists\_ge\_le\_exp\_horizon}: a version of the pre-optimization Freedman
bound (Lemma~\ref{lem:lil_freedman}) whose increment bound $|\Delta M_i|\le c$ is required only for
$i<N$, and whose event is restricted to $m\le N$. It is proved by applying
Lemma~\ref{lem:lil_freedman} to the deterministically stopped process
$M^N_m:=M_{\min(m,N)}$ (Mathlib's \texttt{stoppedProcess} at the constant time $N$), which is a martingale
(\texttt{martingale\_stoppedProcess\_const}), whose increments vanish past $N$ hence are globally bounded by
$c$, and whose predictable quadratic variation agrees with $\langle M\rangle$ below the horizon
(\texttt{predQuadVar\_stoppedProcess\_const\_of\_le}). Here it is used with $N=2^k$, $c=c_k=2\sqrt{2^k}$, the
constrained $\theta_k=1/c_k$, and the increment bound $|\Delta\tilde M_i|\le2\sqrt i\le c_k$ for
$i<2^k$ (\texttt{abs\_truncRespMart\_increment\_le}).
\end{proof}

\begin{lemma}[Block probabilities are summable]\label{lem:trunc_block_summable}
\lean{AlphaRAR.summable_exp_neg_mul_sqrt_add}\leanok
For every $C>0$, $\sum_k \exp(-\tfrac C2\sqrt k+\tfrac{\sigma^2}4)<\infty$.
\end{lemma}
\begin{proof}\leanok
$\exp(-\tfrac C2\sqrt k)$ is summable in $k$: since $\log=o(\sqrt\cdot)$, eventually
$2\log k\le \tfrac C2\sqrt k$, so the terms are $\le k^{-2}$.
\end{proof}

\begin{lemma}[LIL for the truncated martingale]\label{lem:trunc_mart_lil}
\lean{AlphaRAR.ae_eventually_truncRespMart_lt_block,
AlphaRAR.ae_eventually_truncRespMart_le_sqrt_nat_mul_log,
AlphaRAR.ae_eventually_abs_truncRespMart_le_sqrt_nat_mul_log}\leanok
\uses{lem:trunc_block, lem:trunc_block_summable, lem:trunc_qv_le}
$\tilde M_n=O(\sqrt{n\log n})$ a.s.

\textbf{Reusable tool.} The block argument is formalized \emph{abstractly}, for any $L^2$-martingale
$M$ with a square-root-growing increment bound $|\Delta M_i|\le a\sqrt i$ and a linear quadratic
variation bound $\langle M\rangle_n\le v\,n$: the general lemmas
\texttt{measure\_exists\_ge\_le\_exp\_block} (per-block Freedman),
\texttt{ae\_eventually\_lt\_block\_of\_growing} (block Borel--Cantelli),
\texttt{ae\_eventually\_le\_sqrt\_nat\_mul\_log\_of\_growing} (one-sided $O(\sqrt{n\log n})$) and
\texttt{ae\_eventually\_abs\_le\_sqrt\_nat\_mul\_log\_of\_growing} (\emph{two-sided}, via $\pm M$ and
$\langle -M\rangle=\langle M\rangle$) live in \texttt{LILTruncation.lean} and belong upstream. The
truncated martingale $\tilde M$ is the instance $a=2$, $v=V_k$; in particular the two-sided form
\texttt{ae\_eventually\_abs\_truncRespMart\_le\_sqrt\_nat\_mul\_log} ($|\tilde M_n|\le
C\sqrt{n\log n}$) is free.
\end{lemma}
\begin{proof}\leanok
\uses{lem:trunc_block, lem:trunc_block_summable, lem:trunc_qv_le}
Borel--Cantelli over the blocks (Lemmas~\ref{lem:trunc_block}--\ref{lem:trunc_block_summable})
gives that a.s.\ for all large $k$ and every $m\le 2^k$, $\langle\tilde M\rangle_m\le\sigma^2 2^k
\Rightarrow \tilde M_m<\lambda_k$ (formalized as
\texttt{ae\_eventually\_truncRespMart\_lt\_block}, the summable first Borel--Cantelli step). To
bound a fixed time $n$, take the \emph{least} $k$ with $n\le 2^k$; then $2^k\le 2n$ and
$k\le\log_2 n+1$, and the quadratic-variation bound $\langle\tilde M\rangle_n\le V_k N_{n,k}\le V_k
n\le V_k 2^k$ (Lemma~\ref{lem:trunc_qv_le}) supplies the block hypothesis at $m=n\le 2^k$, giving
$\tilde M_n<\lambda_k=C\sqrt{2^k k}\le C'\sqrt{n\log n}$.

\textbf{Scale.} Because the increment bound $|\Delta\tilde M_i|\le2\sqrt i$ \emph{grows}, the block
horizon $m\le 2^k$ forces the least admissible $k$ to satisfy $2^k\ge n$, so time-blocking delivers
the \emph{time} scale $O(\sqrt{n\log n})$ rather than the variance scale
$O(\sqrt{\langle\tilde M\rangle_n\log\langle\tilde M\rangle_n})=O(\sqrt{N_n\log N_n})$ of the
displayed statement. The two coincide when arm $k$ is sampled a positive fraction of the time
($N_{n,k}\asymp n$, the regime of interest); the sharper variance-scale bound would require blocking
by the quadratic variation rather than by the natural time, and is not formalized. The time-scale
form is exactly parallel to the bounded-increment Corollary~\ref{cor:lil_upper_nat}
(\texttt{ae\_eventually\_le\_sqrt\_nat\_mul\_log}).
\end{proof}

\subsection{Assembling}

\begin{lemma}[Truncation to finite variance]\label{lem:lil_truncation}
\lean{AlphaRAR.ae_eventually_abs_respMart_le_sqrt_nat_mul_log,
AlphaRAR.ae_eventually_respMart_le_sqrt_nat_mul_log}\leanok
\uses{lem:trunc_tail_const, lem:trunc_drift, lem:trunc_mart_lil}
$M_n = O(\sqrt{n\log n})$ a.s. The formalized statement is \emph{two-sided},
\texttt{ae\_eventually\_abs\_respMart\_le\_sqrt\_nat\_mul\_log} ($|M_n|\le C\sqrt{n\log n}$
eventually, a.s.), with the one-sided form its immediate corollary. The two-sidedness comes for free
from the reusable growing-increment LIL tool (see Lemma~\ref{lem:trunc_mart_lil}).
\end{lemma}
\begin{proof}\leanok
\uses{lem:trunc_tail_const, lem:trunc_drift, lem:trunc_mart_lil}
By Definition~\ref{def:lil_trunc}, $M_n=\tilde M_n+R_n+\mathrm{Dr}_n$. The main term is
$O(\sqrt{n\log n})$ (Lemma~\ref{lem:trunc_mart_lil}); the tail is $O(1)$
(Lemma~\ref{lem:trunc_tail_const}); the drift is $O(\sqrt{n})$ (Lemma~\ref{lem:trunc_drift}).
Both remainders are $o(\sqrt{n\log n})$, so $M_n=O(\sqrt{n\log n})$. In the regime of interest
$N_{n,k}\asymp n$, this is the $O(\sqrt{N_n\log N_n})=O(\sqrt{\langle M\rangle_n\log\langle
M\rangle_n})$ rate (using $\langle M\rangle_n=\sigma^2 N_n$) required downstream.
\end{proof}

\textbf{Remark.} Lemma~\ref{lem:lil_truncation} is exactly the shape used in
Lemma~\ref{lem:theta_LIL}: there $M = Q_{\cdot,k}$, $D_i = X_{i,k}$,
$\xi_i = \xi_{i,k}-\theta_k$, and $\langle Q_{\cdot,k}\rangle_n = V_k N_{n,k}$, so
the conclusion reads $Q_{n,k} = O(\sqrt{N_{n,k}\log N_{n,k}})$ a.s. The same
statement applied to the bounded martingale $M_{\cdot,k}$
(Definition~\ref{def:M}) covers the use in Lemma~\ref{lem:U_increment_bound}; that direction is
already formalized as Corollary~\ref{cor:lil_upper_nat}.

\textbf{Formalization status.} All of \S\ref{sec:pre_lil_trunc} is now formalized: the tail
(Lemmas~\ref{lem:trunc_tail_summable}--\ref{lem:trunc_tail_const}), the drift
(Lemmas~\ref{lem:trunc_mean_bound}--\ref{lem:trunc_drift}), the martingale and its quadratic
variation (Lemmas~\ref{lem:trunc_mart}--\ref{lem:trunc_qv_le}), the core growing-increment LIL
(Lemmas~\ref{lem:trunc_block}--\ref{lem:trunc_mart_lil}), and the assembling
Lemma~\ref{lem:lil_truncation} (\texttt{ae\_eventually\_respMart\_le\_sqrt\_nat\_mul\_log}). The core
reuses \texttt{measure\_exists\_ge\_le\_exp} on a \emph{deterministically stopped} martingale --- the
one genuinely new device, packaged as \texttt{measure\_exists\_ge\_le\_exp\_horizon} --- followed by
the Borel--Cantelli/repackaging pattern of \S\ref{sec:pre_lil_upper}. The whole chapter delivers the
time-scale $O(\sqrt{n\log n})$ form (see the scale remark in Lemma~\ref{lem:trunc_mart_lil});
the variance-scale $O(\sqrt{\langle M\rangle_n\log\langle M\rangle_n})$, which agrees in the regime
$N_{n,k}\asymp n$, would require blocking by the quadratic variation and is not formalized.
